% !TeX root = ../main.tex
\section{Reduction to the Nernst--Planck--Darcy system of Ignatova and Shu}
\label{app:nernst_planck_darcy_audit}

This appendix verifies the correspondence between
\eqref{eq:nernst_planck_darcy_limit} and the system of
\cite[eqs.~(1.1)--(1.5)]{ignatova2021charged_fluids}.  The specialization of
\sref{nernst_planck_darcy_limit} is retained: one isothermal fluid phase in a
fixed solid frame, vanishing reaction, phase-transformation, and
mechanical-dispersion rates, negligible inertia, gravity, and capillarity,
constant \(\phi_f\), a linear isotropic dielectric response
\(\mbf d_f=\epsilon_f\mbf E\) with constant \(\epsilon_f\), all free charge
assigned to the fluid components, and the electroquasistatic relation
\(\mbf E=-\nabla_{\mbf x}\varphi\).  We further set \(\phi_f=1\),
nondimensionalize the coefficients, take two ionic components with opposite
valences and equal diagonal diffusivities, and adopt the ideal-dilute choice
of the component potentials.  The notation of
\cite{ignatova2021charged_fluids} is related to the present notation by
%
\begin{equation}
	u=\mbf v_f,
	\qquad
	p=\bar p_f,
	\qquad
	\Phi=\varphi,
	\qquad
	c_i=\rho_f^i,
	\qquad
	z_i=z_f^i,
	\qquad
	\rho=\varrho,
	\qquad
	\varrho=\sum_\alpha z_f^\alpha\rho_f^\alpha,
	\qquad
	\varepsilon=\epsilon_f,
	\label{eq:nernst_planck_darcy_notation_identification}
\end{equation}
%
where the concentration \(c_i\), the valence \(z_i\), and the charge density
\(\varrho\) are the nondimensionalized forms of the partial density
\(\rho_f^i\), the charge-to-mass ratio \(z_f^i\), and the fluid charge sum,
respectively, the parameter \(\varepsilon\) of
\cite{ignatova2021charged_fluids}, proportional to the square of the Debye
length, is the nondimensionalized dielectric permittivity \(\epsilon_f\),
and the isotropic diffusion tensor \(\mbf D_{f,{\rm diff}}^{\alpha\beta}\)
is diagonal with common entry \(D_i=D\).

\subsection{Darcy velocity and incompressibility}

Setting \(\phi_f=1\) in \eqref{eq:nernst_planck_darcy_velocity} gives
%
\begin{equation}
	\mbf v_f
	=
	-\frac{\bs\kappa_f}{\mu_f^{\mathrm{visc}}}
	\left(
		\nabla_{\mbf x}\bar p_f-\varrho\mbf E
	\right),
	\qquad
	\varrho=\sum_\alpha z_f^\alpha\rho_f^\alpha .
	\label{eq:nernst_planck_darcy_audit_velocity_rescaled}
\end{equation}
%
The mobility is isotropic and, as in
\cite{ignatova2021charged_fluids}, the friction coefficient is set to one, so
the mobility factor reduces to unity.  The electroquasistatic relation
\(\mbf E=-\nabla_{\mbf x}\varphi\) then gives
%
\begin{equation}
	\mbf v_f
	=
	-\nabla_{\mbf x}\bar p_f-\varrho\nabla_{\mbf x}\varphi .
	\label{eq:nernst_planck_darcy_audit_darcy_law}
\end{equation}
%
With the identifications in
\eqref{eq:nernst_planck_darcy_notation_identification}, equation
\eqref{eq:nernst_planck_darcy_audit_darcy_law} becomes
%
\begin{equation}
	u+\nabla p=-\varrho\nabla\Phi,
	\label{eq:nernst_planck_darcy_audit_recovered_darcy}
\end{equation}
%
which is equation~(1.4) of \cite{ignatova2021charged_fluids}; its charge
density \(\rho\) is written \(\varrho\) here.  The total
mass balance of the fluid phase, obtained by summing
\eqref{eq:nernst_planck_darcy_component_balance} over \(\alpha\) and using
the zero-sum relative fluxes, is
\(\partial\rho_f/\partial t+\nabla_{\mbf x}\cdot(\rho_f\mbf v_f)=0\).
The incompressible constant-density specialization gives
\(\nabla_{\mbf x}\cdot\mbf v_f=0\), which is equation~(1.5) of
\cite{ignatova2021charged_fluids}.

\subsection{Poisson equation and charge density}

Setting \(\phi_f=1\), \(\mbf d_f=\epsilon_f\mbf E\), and
\(\mbf E=-\nabla_{\mbf x}\varphi\) in
\eqref{eq:nernst_planck_darcy_gauss_law} gives
%
\begin{equation}
	-\epsilon_f\Delta_{\mbf x}\varphi
	=
	\varrho,
	\qquad
	\varrho=\sum_\alpha z_f^\alpha\rho_f^\alpha,
	\qquad
	\Delta_{\mbf x}=\nabla_{\mbf x}\cdot\nabla_{\mbf x},
	\label{eq:nernst_planck_darcy_audit_poisson}
\end{equation}
%
because \(\epsilon_f\) is constant.  Nondimensionalization, with the
parameter \(\varepsilon\) of \cite{ignatova2021charged_fluids} identified
with \(\epsilon_f\) as in
\eqref{eq:nernst_planck_darcy_notation_identification}, gives
%
\begin{equation}
	-\varepsilon\Delta\Phi=\varrho,
	\qquad
	\varrho=\sum_i z_i c_i,
	\label{eq:nernst_planck_darcy_audit_recovered_poisson}
\end{equation}
%
which are equations~(1.2)--(1.3) of \cite{ignatova2021charged_fluids}; their
charge density \(\rho\) is written \(\varrho\) here.

\subsection{Nernst--Planck component balances}

For two ionic components, \(N=2\) and the reference component is \(N=2\).
The isotropic, diagonal, equal-diffusivity closure leaves a single term in
the sum over \(\beta\ne N\), so the independent balance
\eqref{eq:nernst_planck_darcy_component_balance} for \(\alpha=1\) becomes
%
\begin{equation}
	\frac{\partial\rho_f^1}{\partial t}
	+
	\nabla_{\mbf x}\cdot
	\left[
		\rho_f^1\mbf v_f
		-
		D\nabla_{\mbf x}
		\left(
			\hat\mu_f^1-\hat\mu_f^2
			+
			\left(z_f^1-z_f^2\right)\varphi
		\right)
	\right]
	=0 .
	\label{eq:nernst_planck_darcy_audit_two_ion_balance}
\end{equation}
%
The ideal-dilute choice
\(\hat\mu_f^\alpha=\hat\mu_f^{\circ\alpha}+\ln\rho_f^\alpha\), with the
constant offset removed by the gradient, separates the electrochemical force
into concentration-diffusion and electromigration parts,
%
\begin{equation}
	\nabla_{\mbf x}
	\left(
		\hat\mu_f^1-\hat\mu_f^2
		+
		\left(z_f^1-z_f^2\right)\varphi
	\right)
	=
	\nabla_{\mbf x}\ln\frac{\rho_f^1}{\rho_f^2}
	+
	\left(z_f^1-z_f^2\right)\nabla_{\mbf x}\varphi .
	\label{eq:nernst_planck_darcy_audit_force_split}
\end{equation}
%
The first term is the concentration-diffusion (Fick) part and the second is
the electromigration part.  Under the opposite-valences specialization
\(z_f^2=-z_f^1\), the electromigration coefficient is
\(z_f^1-z_f^2=2z_f^1\).  The reference-component balance is reconstructed
from the zero-sum closure and the total mass balance, and the per-species
Nernst--Planck form used by \cite{ignatova2021charged_fluids} is
%
\begin{equation}
	\frac{\partial c_i}{\partial t}
	+
	\nabla\cdot
	\left(
		u c_i-D_i\nabla c_i-z_iD_i c_i\nabla\Phi
	\right)
	=0,
	\qquad
	i=1,2,
	\label{eq:nernst_planck_darcy_audit_recovered_np}
\end{equation}
%
which is equation~(1.1) of \cite{ignatova2021charged_fluids}.  The
reference-difference driving force in
\eqref{eq:nernst_planck_darcy_audit_force_split} and the per-species flux in
\eqref{eq:nernst_planck_darcy_audit_recovered_np} are connected by the
mass-averaged relative-flux convention used for the resolved phase velocity:
the per-species electrochemical-potential gradient equals the
reference-difference force plus the common drift carried by the velocity.
Both conventions therefore describe the same coupled two-ion
electrodiffusion system.

In summary, \eqref{eq:nernst_planck_darcy_audit_recovered_np},
\eqref{eq:nernst_planck_darcy_audit_recovered_poisson}, and
\eqref{eq:nernst_planck_darcy_audit_recovered_darcy} reproduce
equations~(1.1), (1.2)--(1.3), and (1.4)--(1.5) of
\cite{ignatova2021charged_fluids}, respectively.
